\documentclass[10pt,a4paper]{amsart}
\usepackage[margin=19mm]{geometry}
\usepackage{amsmath,amssymb,mathtools}
\usepackage[scr=rsfs,cal=euler]{mathalfa}
\usepackage{xfrac}
\usepackage[unicode,hidelinks,bookmarks=false]{hyperref}
\newcommand{\ie}{i.e.\ }
\newcommand{\cf}{cf.\ }
\newcommand{\resp}{respectively\ }
\newcommand{\Subsection}{\subsection}
\providecommand{\nobreakdash}{\nobreak\hbox{-}}
\setlength{\emergencystretch}{2em}
\multlinegap0pt
\usepackage{bm}
\usepackage{bbm}
\usepackage{dsfont}
\usepackage{xspace}
\usepackage{cancel}
\usepackage{mathtools}
\usepackage{amsthm}
\makeatletter
\@ifundefined{thm}{\newtheorem{thm}{Thm}}{}
\@ifundefined{lem}{\newtheorem{lem}[thm]{Lemma}}{}
\@ifundefined{prop}{\newtheorem{prop}[thm]{Proposition}}{}
\makeatother
\makeatletter
\expandafter\ifx\csname enumerit\endcsname\relax
\newenvironment{enumerit}{\begin{list}{{\hfill\rm(\roman{cnt})\hfill}}{%
\settowidth{\labelwidth}{{\rm(iv)}}\leftmargin=\labelwidth%
\advance\leftmargin by \labelsep\rightmargin=0pt\usecounter{cnt}}}{\end{list}} \makeatletter
\fi
\expandafter\ifx\csname pf\endcsname\relax
\newenvironment{pf}{\proof}{\endproof}
\fi
\makeatother
\title[Zero-sum Inverse Realization and Property~(P) under Join Ope]{Zero-sum Inverse Realization and Property~(P) under Join Operations}
\author{G. Arunkumar, Anubhab Pahari,, Puja Samanta}
\date{Source: arXiv:2608.10661v1 (CC BY 4.0)}
\begin{document}
\maketitle

\noindent\emph{Source and licence.} Zero-sum Inverse Realization and Property~(P) under Join Operations. Version arXiv:2608.10661v1 (CC BY 4.0), distributed under the Creative Commons Attribution 4.0 International licence, as recorded in the arXiv licence metadata for this exact version and shown on the abstract page https://arxiv.org/abs/2608.10661v1. Full rights notice: licenses/arxiv-2608.10661-CC-BY-4.0.txt.\par\smallskip

\noindent\textit{Editorial scope.} Counted unit: the Prop whose source label is \texttt{prop: nonzero isotropic vector} in the article (source0.tex lines 501--504, source label \texttt{prop: nonzero isotropic vector}), delivered together with the article's own complete original proof of it (source0.tex lines 506--696, one \texttt{proof} environment of 191 lines). No mathematics is written, repaired or re-derived: statement and proof are the article's own text line for line. Required context retained uncounted: avoidance (lines 442--444). Every definition, display and label of those lines is retained unchanged. Omitted from this package and not redistributed: 1--441, 445--500, 505--505, 697--1556. Nothing inside a retained excerpt is omitted or altered: every retained body is delivered exactly as the article's lines are written, so no omission receipt of any kind accompanies this package. Citations: the retained excerpts carry no citation command at all, so no bibliography entry of the article is transcribed and no reference is redistributed. Reference adaptation: none was needed and none was made; every reference inside the delivered unit resolves inside it, so no unresolved reference or citation is produced. Printed slips kept literal: no printed word, symbol, hypothesis or number of the counted statement or of its proof was altered. Layout and class: the article's own class is \texttt{amsart} with packages including amsthm, booktabs, amssymb, graphics, tikz, latexsym, multicol, verbatim, enumerate, accents, cite, enumitem, array, mathtools; the wrapper is the lane's pinned \texttt{amsart} editorial wrapper at 10pt on A4, which restates the theorem-like environments the retained text needs and adds the article's own macro definitions that the retained text needs (none), with extra wrapper packages (bm, bbm, dsfont, xspace, cancel, mathtools). The delivered numbering is the unit's own. Assets: no retained span contains a figure environment, a graphics-inclusion command, a PDF-page inclusion, a table or a TikZ picture; no image file is copied into this package, the package declares no asset dependency and no image file is redistributed with it. Licence: version arXiv:2608.10661v1 of this article is distributed under the Creative Commons Attribution 4.0 International licence, as recorded in the arXiv licence metadata for this exact version and shown on the abstract page https://arxiv.org/abs/2608.10661v1. A full rights notice is delivered at licenses/arxiv-2608.10661-CC-BY-4.0.txt. Characters: every retained excerpt is pure ASCII, so the delivered engine, XeLaTeX with its default Latin Modern text font, reports no missing character.
\begin{lem}\label{avoidance}
For $d\ge 1$, let $f\in \mathbb{R}[x_1,\dots,x_d]$ be a nonzero polynomial. Then $f$ does not vanish on any nonempty open subset of $\mathbb{R}^d$, where $\mathbb{R}^d$ is equipped with the standard Euclidean topology.
\end{lem}

\begin{prop}\label{prop: nonzero isotropic vector}
Let $Q$ be a nondegenerate indefinite quadratic form on $\mathbb{R}^n$ with
$n\ge3$. Then $Q$ admits a fully supported isotropic vector $x \in \mathbb{R}^n$.
\end{prop}

\begin{pf}
Write $$x=\begin{pmatrix}x'\\x_n\end{pmatrix}\in\mathbb{R}^n,\  \text{where}
\
x'=(x_1,\dots,x_{n-1})^\top\in\mathbb{R}^{n-1}.
$$
Let $B$ denote the symmetric bilinear form associated with $Q$, and define
\[
\phi(x')=Q\begin{pmatrix}x'\\0\end{pmatrix},\qquad
\ell(x')=B\left(e_n,\begin{pmatrix}x'\\0\end{pmatrix}\right),\qquad
c=Q(e_n).
\]
Then
\begin{equation}\label{eq:Q-decomposition}
Q\begin{pmatrix}x'\\x_n\end{pmatrix}
=
\phi(x')+2\ell(x')x_n+cx_n^2.
\end{equation}

We first show that $\phi\not\equiv 0$. Suppose, to the contrary, that
$\phi\equiv 0$. Then $Q$ vanishes on the hyperplane
\(
H=\operatorname{span}\{e_1,\dots,e_{n-1}\}.
\)
If $v,w\in H$, then
\[
Q(v+w)=Q(v)+Q(w)+2B(v,w).
\]
Since $Q$ vanishes on $H$, we get $B(v,w)=0$ for all $v,w\in H$. Hence
\[
H\subseteq H^\perp,
\]
where
\(
H^\perp=\{x\in\mathbb{R}^n:B(x,h)=0\text{ for all }h\in H\}
\)
denotes the orthogonal complement of $H$ with respect to $B$.
Since $Q$ is nondegenerate, the associated bilinear form $B$ is nondegenerate. Hence
\[
\dim H^\perp=n-\dim H=1.
\]
Thus
\[
n-1=\dim H\le \dim H^\perp=1,
\]
which is impossible because $n\ge 3$. Therefore, $\phi$ is a nonzero polynomial.

We now consider two cases.

\medskip

\noindent\textbf{Case 1: $c\neq 0$.}
In this case, equation \eqref{eq:Q-decomposition} can be rewritten as
\[
Q\begin{pmatrix}x'\\x_n\end{pmatrix}
=
c\left(x_n+\frac{\ell(x')}{c}\right)^2
-\frac{1}{c}\Delta(x'),
\]
where
\[
\Delta(x')=\ell(x')^2-c\phi(x').
\]

We next show that $\Delta(x')>0$ for some
$x'\in\mathbb R^{n-1}$. Since $Q$ is indefinite, it
attains a value of sign opposite to $c$. If $c>0$, choose
$u=(u',u_n)^\top\in\mathbb{R}^n$ such that $Q(u)<0$. Then
\[
Q(u)=
c\left(u_n+\frac{\ell(u')}{c}\right)^2
-\frac{1}{c}\Delta(u').
\]
Since the first term on the right-hand side is nonnegative, we must have
\(
-\frac{1}{c}\Delta(u')<0.
\)
As $c>0$, this gives $\Delta(u')>0$.
Similarly, if $c<0$, choose $u\in\mathbb{R}^n$ such that $Q(u)>0$. Then
\(
c\left(u_n+\frac{\ell(u')}{c}\right)^2\le 0,
\)
and hence the equality above forces $\Delta(u')>0$.

Therefore
\(
U=\{x'\in\mathbb{R}^{n-1}:\Delta(x')>0\}
\)
is a nonempty open subset of $\mathbb{R}^{n-1}$.
Consider the polynomial
\(
F(x')=\phi(x')x_1\cdots x_{n-1}.
\)
Since $\phi\not\equiv 0$, the polynomial $F$ is nonzero. By
Lemma~\ref{avoidance}, $F$ cannot vanish on the nonempty open set $U$.
Therefore, there exists
\(
x^\ast=(x_1^\ast,\dots,x_{n-1}^\ast)^\top\in U
\)
such that
\(
F(x^\ast)\neq 0.
\)
Thus
\[
\phi(x^\ast)\neq 0
\quad\text{and}\quad
x_i^\ast\neq 0
\quad\text{for all } i=1,\dots,n-1.
\]

%For this choice of $x^\ast$, consider the equation {\color{blue}need to write more}
Fix \(x'=x^\ast\). By \eqref{eq:Q-decomposition}, treating the last coordinate as a variable
\(t\in\mathbb{R}\), we obtain the quadratic equation
\[
ct^2+2\ell(x^\ast)t+\phi(x^\ast)=0.
\]
Its discriminant is
\[
4\ell(x^\ast)^2-4c\phi(x^\ast)=4\Delta(x^\ast)>0.
\]
Hence it has two real roots. Moreover, since $\phi(x^\ast)\neq 0$, neither
root is zero. Choosing $x_n$ to be either root gives
\[
Q\begin{pmatrix}x^\ast\\x_n\end{pmatrix}=0,
\]
and all coordinates of this vector are nonzero.

\medskip

\noindent\textbf{Case 2: $c=0$.}
In this case, equation \eqref{eq:Q-decomposition} reduces to
\[
Q\begin{pmatrix}x'\\x_n\end{pmatrix}
=
\phi(x')+2\ell(x')x_n.
\]

We first show that $\ell$ is not identically zero. Indeed,
\[
\ell(x')=
B\left(e_n,\begin{pmatrix}x'\\0\end{pmatrix}\right)
=
\sum_{j=1}^{n-1}B(e_n,e_j)x_j.
\]
If $\ell\equiv 0$, then
\[
B(e_n,e_j)=0
\quad\text{for all } 1\le j\le n-1.
\]
Since $c=0$, we also have
\[
B(e_n,e_n)=Q(e_n)=0.
\]
Therefore,
\[
B(e_n,e_j)=0
\quad\text{for all } 1\le j\le n.
\]
Thus, $e_n$ is orthogonal to all vectors in $\mathbb{R}^n$, which means that
$e_n$ lies in the radical of $B$. This contradicts the nondegeneracy of $Q$.
Hence $\ell\not\equiv 0$.

Consider the polynomial
\[
G(x')=\phi(x')\ell(x')x_1\cdots x_{n-1}.
\]
Observe that $G\not\equiv 0$. By Lemma~\ref{avoidance}, there exists
$x^\ast\in\mathbb{R}^{n-1}$ such that
\[
\phi(x^\ast)\neq 0,\qquad
\ell(x^\ast)\neq 0,\qquad
x_1^\ast\cdots x_{n-1}^\ast\neq 0.
\]
Set
\[
x_n=-\frac{\phi(x^\ast)}{2\ell(x^\ast)}.
\]
Then $x_n\neq 0$, and
\[
Q\begin{pmatrix}x^\ast\\x_n\end{pmatrix}
=
\phi(x^\ast)+2\ell(x^\ast)
\left(-\frac{\phi(x^\ast)}{2\ell(x^\ast)}\right)
=0.
\]
Therefore,
\(
\begin{pmatrix}x^\ast\\x_n\end{pmatrix}
\)
is an isotropic vector for $Q$ whose coordinates are all nonzero.
\end{pf}

\end{document}
